\section*{Hoofdstuk \ref{ch:g12:metamorphic-rocks} --- Metamorfe gesteenten}

\begin{solution}{exo:g12:metamorphic-rocks:1}
(a) \omterm{def:g12:metamorphic-rocks:polymorph}{polymorfen}; (b) een \omterm{def:g12:metamorphic-rocks:index-mineral}{isograad}; (c) een \omterm{def:g12:metamorphic-rocks:contact-regional}{aureool}; (d) een \omterm{def:g12:metamorphic-rocks:foliation}{foliatie};
(e) de \omterm{def:g12:metamorphic-rocks:lithostatic}{lithostatische druk}.
\end{solution}

\begin{solution}{exo:g12:metamorphic-rocks:2}
(1) marmer; (2) kwartsiet; (3) leisteen; (4) gneis; (5) hoornfels.
\end{solution}

\begin{solution}{exo:g12:metamorphic-rocks:3}
$P = \rho g z$: op \qty{10}{km} $2830 \times 9.81 \times \num{10000} \approx
\qty{2.78e8}{Pa} = \qty{0.278}{GPa}$; op \qty{20}{km} twee keer zoveel,
\qty{0.555}{GPa}.
\end{solution}

\begin{solution}{exo:g12:metamorphic-rocks:4}
(1) Kyaniet (de grens tussen kyaniet en andalusiet ligt bij \qty{400}{\celsius} rond \qty{0.25}{GPa}).
(2) Andalusiet. (3) Sillimaniet.
\end{solution}

\begin{solution}{exo:g12:metamorphic-rocks:5}
(1) Leisteen, fylliet, schist, gneis. (2) In de schist. (3) De gneis.
\end{solution}

\begin{solution}{exo:g12:metamorphic-rocks:6}
(1) $2830 \times 9.81 \times \num{35000} \approx \qty{9.72e8}{Pa}$, dat is
$\num{9.72e8}/101325 \approx \num{9600}$ atmosfeer. (2) $1025 \times 9.81
\times 4000 \approx \qty{4.02e7}{Pa} = \qty{40.2}{MPa}$. (3) $972/40.2
\approx 24$ keer.
\end{solution}

\begin{solution}{exo:g12:metamorphic-rocks:7}
(1) $z = \num{0.375e9}/(2830 \times 9.81) \approx \qty{13.5}{km}$.
(2) $(504 - 10)/13.5 \approx \qty{36.6}{\celsius}$ per kilometer. (3) Nee:
het gesteente zou precies in één punt van het diagram moeten liggen; zulke gesteenten
zijn zeldzaam.
\end{solution}

\begin{solution}{exo:g12:metamorphic-rocks:8}
(1) $23/3 \approx \qty{8}{mm}$. (2) Ongeveer 8 miljoen jaar. (3) Het gesteente
bleef miljoenen jaren op diepte, bij een hoge temperatuur.
\end{solution}

\begin{solution}{exo:g12:metamorphic-rocks:9}
(1) Vier: de isograden van granaat, stauroliet, kyaniet en sillimaniet, in die
volgorde. (2) Naar het westen. (3) Aan het westelijke uiteinde, in de sillimanietzone. Zijn
gesteenten overschreden de grens tussen kyaniet en sillimaniet: begraven op ongeveer \qty{20}{km}
(\qty{0.56}{GPa}) waren ze heter dan ongeveer \qty{590}{\celsius}.
\end{solution}

\begin{solution}{exo:g12:metamorphic-rocks:10}
(1) Modder wordt hoornfels, kalksteen wordt marmer, zandsteen wordt
kwartsiet. (2) Ze werden verhit zonder te worden samengeperst: geen gerichte druk
legde hun mineralen op één lijn. (3) De gesteenten eromheen zijn al heet, dus
verhit het magma ze maar weinig meer.
\end{solution}

\begin{solution}{exo:g12:metamorphic-rocks:11}
(1) Dwars op de duw: hun platte vlakken zijn naar het oosten en het westen gericht, zodat de schilfers
in noord--zuidvlakken staan. (2) Verticaal, noord--zuid lopend.
(3) Korrels calciet en kwarts zijn plat noch langwerpig, dus zelfs wanneer ze
onder samendrukking herkristalliseren, is er niets om op één lijn te leggen.
\end{solution}

\begin{solution}{exo:g12:metamorphic-rocks:12}
(1) $z = (T - 10)/30$, dus $P = 0.0278\,(T - 10)/30 \approx
0.000927\,(T - 10)$. (2) $0.000927\,(T - 10) = 0.577 + 0.00209\,(T - 600)$
geeft $0.001163\,T \approx 0.668$, dus $T \approx \qty{574}{\celsius}$, $P
\approx 0.000927 \times 564 \approx \qty{0.523}{GPa}$ en $z = (574 - 10)/30
\approx \qty{18.8}{km}$. (3) Kyaniet boven die diepte, sillimaniet eronder.
\end{solution}

\begin{solution}{exo:g12:metamorphic-rocks:13}
(1) $10 + 10 \times 35 = \qty{360}{\celsius}$. (2) $10 + 30 \times 35 =
\qty{1060}{\celsius}$. (3) De druk op \qty{35}{km} is ongeveer
\qty{0.97}{GPa}; bij \qty{360}{\celsius} ligt dat punt diep in het
kyanietveld. Het tweede punt ligt ver rechts van het diagram, in het
sillimanietveld, bij temperaturen waarbij korstgesteenten beginnen te smelten.
(4) Ze hebben een hoge druk bij een lage temperatuur nodig, een combinatie die alleen
de koude korst levert die in een \omterm{def:g11:subduction-zones:subduction}{subductiezone} omlaag wordt gevoerd.
\end{solution}

\begin{solution}{exo:g12:metamorphic-rocks:14}
(1) Kyaniet $1000/3.61 \approx \qty{277}{cm^3}$; sillimaniet $1000/3.24
\approx \qty{309}{cm^3}$; andalusiet $1000/3.15 \approx \qty{317}{cm^3}$.
(2) $(317 - 277)/317 \approx 0.127$, ongeveer \qty{13}{\percent}. (3) Bij hoge
druk wordt de compactste rangschikking van de atomen begunstigd; kyaniet
pakt dezelfde atomen in de minste ruimte.
\end{solution}

\begin{solution}{exo:g12:metamorphic-rocks:15}
(1) \qty{0.555}{GPa} op \qty{20}{km}, \qty{0.694}{GPa} op \qty{25}{km}.
(2) $T = 600 + (P - 0.577)/0.00209$: ongeveer \qty{590}{\celsius} bij
\qty{0.555}{GPa} en \qty{656}{\celsius} bij \qty{0.694}{GPa}. Bij
\qty{700}{\celsius} waren de gesteenten heter dan beide, in het
sillimanietveld: dat klopt. (3) $(700 - 10)/20 \approx \qty{34.5}{\celsius}$ en
$(700 - 10)/25 \approx \qty{27.6}{\celsius}$ per kilometer, dicht bij de
gewone gradiënt. (4) $\qty{20}{km}$ tot $\qty{25}{km}$ in 40 miljoen jaar:
0.50 tot \qty{0.63}{km} per miljoen jaar, dat is 0.50 tot \qty{0.63}{mm} per
jaar.
\end{solution}

\begin{solution}{pb:g12:metamorphic-rocks:1}
\textbf{1.} Het is gefolieerd, en zijn mica's zijn groot genoeg om te zien. De
mica's, muscoviet en biotiet, vormen de \omterm{def:g12:metamorphic-rocks:foliation}{schistositeit}.

\textbf{2.} \ce{Al2SiO5}; het zijn \omterm{def:g12:metamorphic-rocks:polymorph}{polymorfen}.

\textbf{3.} De volumes van een kilogram staan in de verhouding $3.15/3.61 \approx
0.873$: kyaniet is ongeveer \qty{13}{\percent} kleiner. Druk begunstigt
kyaniet.

\textbf{4.} De kyanietzone. Op zijn weg naar beneden overschreed het de isograden van biotiet,
granaat, stauroliet en kyaniet.

\textbf{5.} Ongeveer 8 miljoen jaar.

\textbf{6.} $P = 0.577 + 0.00209 \times 0 = \qty{0.577}{GPa}$.

\textbf{7.} \qty{0.70}{GPa} ligt boven \qty{0.577}{GPa}, aan de kyanietkant van
de grens tussen kyaniet en sillimaniet; het ligt ook ver boven het andalusietveld,
dat bij \qty{600}{\celsius} onder ongeveer \qty{0.23}{GPa} ligt.

\textbf{8.} \qty{0.577}{GPa}.

\textbf{9.} $0.70 = 0.577 + 0.00209\,(T - 600)$, dus $T = 600 +
0.123/0.00209 \approx \qty{659}{\celsius}$.

\textbf{10.} \qty{7.0}{kbar}; $\num{7.0e8}/101325 \approx \num{6900}$
atmosfeer.

\textbf{11.} $\rho g = 2830 \times 9.81 \approx \qty{27800}{Pa}$ per meter,
\qty{27.8}{MPa} per kilometer.

\textbf{12.} $z = \num{0.70e9}/\num{27762} \approx \qty{25200}{m}$, ongeveer
\qty{25.2}{km}.

\textbf{13.} $\num{0.577e9}/\num{27762} \approx \qty{20.8}{km}$.

\textbf{14.} $\num{0.70e9}/(2700 \times 9.81) \approx \qty{26.4}{km}$, ongeveer
\qty{5}{\percent} dieper ($26.4/25.2 \approx 1.048$).

\textbf{15.} In de korst: \qty{25}{km} is minder dan zelfs een gewone
korstdikte.

\textbf{16.} $(600 - 10)/25.2 \approx \qty{23.4}{\celsius}$ per kilometer.

\textbf{17.} Ze is veel warmer dan een \omterm{def:g11:subduction-zones:subduction}{subductiezone} (onder
\qty{10}{\celsius} per kilometer) en iets kouder dan een gewone korst
(\qty{30}{\celsius}), ver van een vulkanisch gebied (40 tot \qty{50}{\celsius}): een
gesteente dat diep in de verdikte korst van een gebergtegordel is begraven, door \omterm{def:g12:metamorphic-rocks:contact-regional}{regionale
metamorfose}.

\textbf{18.} $10 + 30 \times 25.2 \approx \qty{766}{\celsius}$. De
grens tussen kyaniet en sillimaniet ligt, doorgetrokken, bij die temperatuur op $0.577 + 0.00209 \times 166
\approx \qty{0.92}{GPa}$; \qty{0.70}{GPa} ligt eronder,
in het sillimanietveld. Een gewone \omterm{def:g11:lithosphere-asthenosphere:geotherm}{geotherm} zou sillimaniet hebben gemaakt.

\textbf{19.} $\qty{25.2}{km} = \qty{2.52e7}{mm}$; bij \qty{0.5}{mm} per jaar
ongeveer 50 miljoen jaar.

\textbf{20.} Terwijl het gesteente opsteeg, koelde het af, dus vertraagden zijn reacties, en
het water dat de omgekeerde reacties nodig hebben, was op de weg naar beneden
verdreven.

\textbf{21.} De kyanietschist kristalliseerde op ongeveer \textbf{\qty{25}{km}
diepte}, bij ongeveer \qty{600}{\celsius}.
\end{solution}
